\section{Introduction}
\label{sec:intro}

Cloud RAN centralizes baseband processing in data-center servers, replacing purpose-built radio hardware with software running on commodity accelerators. The promise is economic: pool processing resources across many cells, scale elastically with demand, and exploit commercial hardware economics. GPUs have emerged as the leading accelerator for this workload. NVIDIA Aerial~\cite{nvidia_aerial} provides a CUDA-accelerated PHY stack that processes 5G NR uplink and downlink channels on GPUs, and production deployments use it to serve multiple cells from a single GPU.

The central question for multi-cell GPU vRAN is scheduling: how should kernels from multiple cells be assigned to GPU resources to meet the strict real-time deadlines of 5G? Each 5G NR slot must complete within 0.5~ms at 120~kHz subcarrier spacing; the uplink equalization pipeline alone must finish in a fraction of that budget. When a GPU serves one cell, the kernel runs in isolation and meets its deadline. When it serves four cells, the natural approach is to run all four concurrently using CUDA streams and rely on the GPU's hardware scheduler to interleave them.

This approach fails because it treats the GPU as a compute-only resource. We profile the Aerial uplink equalization pipeline on an NVIDIA A100 and find that the dominant kernel---the coefficient application step of MMSE equalization---is memory-bound, not compute-bound. The kernel has an arithmetic intensity of 4.32~FLOP/byte, well below the A100's Ridge point of $\sim$161~FLOP/byte for FP16 Tensor Cores. Even after replacing the scalar FP32 inner loop with WMMA FP16 Tensor Core operations---a 5.5$\times$ speedup that reduces latency from 0.146~ms to 0.027~ms---the kernel achieves only 2,681~GFLOP/s, or 0.86\% of the 312~TFLOP/s Tensor Core peak. Profiling with Nsight Compute confirms that the bottleneck is L1 cache throughput (93.6\% of peak in the scalar version, 55.0\% after Tensor Core optimization) and DRAM bandwidth (21.2\% after optimization), not compute pipeline occupancy.

The implication for multi-cell scheduling is direct: if a single cell's equalization kernel already consumes a large fraction of memory bandwidth, concurrent execution of multiple cells will saturate the memory subsystem. Our measurements confirm this: per-cell latency degrades by 1.50$\times$ at 2~cells, 2.43$\times$ at 4~cells, and 3.97$\times$ at 8~cells when all cells run concurrently on separate CUDA streams. SM-based scheduling---which assigns cells to SM partitions or relies on the GPU's warp scheduler to interleave---does not address this because the contention is in the memory hierarchy, not in the SMs.

\textbf{This paper.} We present a systematic experimental study of multi-cell memory contention in GPU vRAN and propose a research line for a memory-hierarchy-aware scheduler. Through six controlled experiments on an A100, we make three key findings that distinguish vRAN memory contention from the well-studied DNN inference case:

\begin{itemize}[leftmargin=*,nosep]
\item \textbf{L2 cache capacity, not HBM bandwidth, is the contention threshold} (\S\ref{sec:motivation}). Two 100~MHz cells with a combined working set of 33~MB (fitting in the 42~MB L2) degrade only 1.12$\times$ when paired with a small 5~MHz cell, but 1.52$\times$ when paired with another 100~MHz cell that overflows L2. This is fundamentally different from DNN serving, where working sets far exceed cache and HBM is the sole bottleneck.

\item \textbf{Pipeline stage ordering reduces contention} (\S\ref{sec:motivation}). The vRAN pipeline has three stages with different working set sizes (Gram: 13.8~MB, Coef: 16.5~MB, LLR: 5.6~MB). Launching the smaller-working-set stage first reduces degradation from 1.50$\times$ (same-stage) to 1.16$\times$ (cross-stage with ordering), because the smaller kernel occupies less L2, leaving capacity for the larger kernel.

\item \textbf{Online contention prediction is feasible without PMUs} (\S\ref{sec:motivation}). Just 5~runtime latency samples predict actual contention with $<$3\% error for $N \geq 2$ cells. This enables adaptive batch sizing: the scheduler starts with a large batch, measures contention, and splits if degradation exceeds a threshold (2.0$\times$).
\end{itemize}

Based on these findings, we propose a scheduler design combining: (1)~L2-aware stage co-scheduling that orders pipeline stages by working set size, (2)~online contention prediction from timing-only measurements, and (3)~adaptive batch sizing with deadline guarantees. We present the experimental evidence, the calibrated contention model, and the development plan toward a full system evaluation (\S\ref{sec:design}--\S\ref{sec:evaluation}).
